% vim:ft=tex
\documentclass{beamer}
\usetheme{metropolis}

\usepackage{annotate-equations}

\usepackage{tikz}
\usepackage{spath3}

% \tikzstyle{dot}=[
% circle, draw=black, fill=black, minimum size=2mm, inner sep=0mm]

% \tikzset{
% 	state/.style={
% 			shape=isosceles triangle,
% 			draw=black,
% 			minimum height=0.75cm,
% 			minimum width=1.0cm,
% 			shape border rotate=#1,
% 			% align=center,
% 			% shift={(0,-0.25cm)},
% 			isosceles
% 			triangle
% 			stretches,
% 			inner
% 			sep=0pt,
% 		},
% }

% macros
% must-haves
\usepackage{amssymb}
\usepackage{amsmath}
\usepackage{amsthm}
\usepackage{amsfonts}
\usepackage{amsxtra}

\usepackage{todonotes}
\usepackage{multicol}
\usepackage{xfrac}

% fonts
\usepackage{bbm}
\usepackage{graphicx}

\definecolor{cyan}{RGB}{102,204,238}
\definecolor{red}{RGB}{238,102,119}
\definecolor{green}{RGB}{34,136,51}

% links/refs
\usepackage[
	unicode=true,
	pdfusetitle,
	bookmarks=true,
	bookmarksnumbered=false,
	breaklinks=false,
	backref=false,
	colorlinks=true,
	linkcolor=blue,
	citecolor=blue,
	urlcolor=blue,
	final
]{hyperref}
\usepackage{bookmark}
\usepackage[
	backend=biber,
	style=alphabetic
]{biblatex}
\addbibresource{references.bib}
\usepackage[capitalise]{cleveref}

% figs
\usepackage{float}
\usepackage{tikz-cd}
\usepackage{caption}
\usepackage{subcaption}
\usepackage{mathpartir}
\usepackage{adjustbox}
\usetikzlibrary{matrix, arrows, cd, patterns}

% diagrams
../writeup/diagrams.tex

% formatting
\usepackage{enumitem}

% crypto
\usepackage[
	n,
	advantage,
	operators,
	sets,
	adversary,
	landau,
	probability,
	notions,
	logic,
	ff,
	mm,
	primitives,
	events,
	complexity,
	oracles,
	asymptotics,
	keys
]{cryptocode}

% cross-reference format
\creflabelformat{equation}{#2#1#3}
\newlist{questions}{enumerate}{1}
\setlist[questions]{label = \arabic*., ref=\arabic*}

\crefname{questionsi}{question}{questions}
\Crefname{questionsi}{Question}{Questions}

% quality of life
\renewcommand{\emptyset}{\varnothing}
\newcommand{\denote}[1]{\llbracket #1 \rrbracket}

% diagrams
\newcommand{\punctuation}[1]{\makebox[0pt][l]{#1}}

% font conventions
\newcommand{\cat}[1]{{\mathcal{#1}}}
\newcommand{\scat}[1]{{\textsc{\text{#1}}}}
\newcommand{\alg}[1]{{\texttt{#1}}}

\DeclareMathOperator{\cod}{cod}
\DeclareMathOperator{\dom}{dom}
\DeclareMathOperator{\ca}{ca}
\DeclareMathOperator{\view}{view}
\DeclareMathOperator{\exec}{exec}
\DeclareMathOperator{\ideal}{ideal}
\let\state\relax
\newcommand{\state}[1]{\operatorname*{st}(#1)}
\newcommand{\onecomb}[1]{\operatorname*{1-comb}(#1)}
\newcommand{\ncomb}[1]{\operatorname*{n-comb}(#1)}
\newcommand{\ncombb}[1]{\operatorname*{n-comb!}(#1)}
\newcommand{\ncomba}[1]{\operatorname*{n-comb}^*(#1)}
\let\prot\relax
\newcommand{\prot}[2]{\operatorname{prot}_{#1}(#2)}
\newcommand{\prota}[2]{\operatorname{prot}^*_{#1}(#2)}
\newcommand{\protb}[2]{\operatorname{prot}!_{#1}(#2)}
\newcommand{\pred}[1]{\operatorname{pred}(#1)}
\newcommand{\oprot}[2]{\overline{\operatorname{prot}}_{#1}(#2)}
\newcommand{\oprota}[2]{\overline{\operatorname{prot}}^*_{#1}(#2)}
\newcommand{\oprotb}[2]{\overline{\operatorname{prot}}!_{#1}(#2)}

% theorems
\theoremstyle{plain}
\newtheorem{thm}{Theorem}[chapter]
\newtheorem{lemma}[thm]{Lemma}
\newtheorem{prop}[thm]{Proposition}
\newtheorem{conj}[thm]{Conjecture}
\newtheorem{claim}[thm]{Claim}
\newtheorem{cor}{Corollary}[thm]

\theoremstyle{definition}
\newtheorem{dfn}[thm]{Definition}
\newtheorem{con}[thm]{Construction}
\newtheorem{ex}[thm]{Example}
\newtheorem{_game}[thm]{Game}
\newtheorem{_algo}[thm]{Algorithm}

\theoremstyle{remark}
\newtheorem{rmk}[thm]{Remark}
\newtheorem*{ntn}{Notation}

% operators
\def\<{\langle}
\def\>{\rangle}
\newcommand{\xcong}[1]{\overset{#1}{\cong}}
\newcommand{\xequiv}[1]{\overset{#1}{\equiv}}
\newcommand{\xequals}[1]{\overset{#1}{=}}
\def\cind{\xequiv{\text{c}}}

\def\into{\hookrightarrow}
\def\onto{\twoheadrightarrow}
\def\from{\leftarrow}
\def\To{\Rightarrow}
\def\From{\Leftarrow}
\newcommand{\xto}[1]{\xrightarrow{#1}}
\newcommand{\xinto}[1]{\xhookrightarrow{#1}}

\DeclareMathOperator{\id}{id}
\DeclareMathOperator{\Hom}{Hom}
\DeclareMathOperator{\pr}{pr}

% crypto stuff
\usepackage{mdframed}
\newenvironment{game}
{\begin{mdframed}\begin{_game}\mbox{}\begin{enumerate}[itemsep=-0.2em]}
				{\end{enumerate}\end{_game}\end{mdframed}}
\newenvironment{algo}
{\begin{mdframed}\begin{_algo}\mbox{}\begin{enumerate}[itemsep=-0.2em]}
				{\end{enumerate}\end{_algo}\end{mdframed}}

% alphabets
\def\cA{\mathcal{A}}
\def\cB{\mathcal{B}}
\def\cC{\mathcal{C}}
\def\cD{\mathcal{D}}
\def\cE{\mathcal{E}}
\def\cF{\mathcal{F}}
\def\cG{\mathcal{G}}
\def\cH{\mathcal{H}}
\def\cI{\mathcal{I}}
\def\cJ{\mathcal{J}}
\def\cK{\mathcal{K}}
\def\cL{\mathcal{L}}
\def\cM{\mathcal{M}}
\def\cN{\mathcal{N}}
\def\cO{\mathcal{O}}
\def\cP{\mathcal{P}}
\def\cQ{\mathcal{Q}}
\def\cR{\mathcal{R}}
\def\cS{\mathcal{S}}
\def\cT{\mathcal{T}}
\def\cU{\mathcal{U}}
\def\cV{\mathcal{V}}
\def\cW{\mathcal{W}}
\def\cX{\mathcal{X}}
\def\cY{\mathcal{Y}}
\def\cZ{\mathcal{Z}}

\def\AA{\mathbb{A}}
\def\BB{\mathbb{B}}
\def\CC{\mathbb{C}}
\def\DD{\mathbb{D}}
\def\EE{\mathbb{E}}
\def\FF{\mathbb{F}}
\def\GG{\mathbb{G}}
\def\HH{\mathbb{H}}
\def\II{\mathbb{I}}
\def\JJ{\mathbb{J}}
\def\KK{\mathbb{K}}
\def\LL{\mathbb{L}}
\def\MM{\mathbb{M}}
\def\NN{\mathbb{N}}
\def\OO{\mathbb{O}}
\def\PP{\mathbb{P}}
\def\QQ{\mathbb{Q}}
\def\RR{\mathbb{R}}
\def\SS{\mathbb{S}}
\def\TT{\mathbb{T}}
\def\UU{\mathbb{U}}
\def\VV{\mathbb{V}}
\def\WW{\mathbb{W}}
\def\XX{\mathbb{X}}
\def\YY{\mathbb{Y}}
\def\ZZ{\mathbb{Z}}

\def\kk{\Bbbk}
\def\ind{\mathbbm{1}}

% itemize spacing
% https://tex.stackexchange.com/a/650378
\let\tempone\itemize
\let\temptwo\enditemize
\renewenvironment{itemize}{
\tempone\setlength{\itemsep}{-1ex}}
{\temptwo}{% endraw %}

%% PROTOCOLS %%
% \createprocedurebox{pro}{center, boxed}{}{}{}
% \newsavebox{\mybox}
% \newenvironment{protocol}[1]
% {\begin{lrbox}{\mybox}\procedure{#1}\begin{minipage}{0.9\textwidth}}
% {\end{minipage}\end{lrbox}\fbox{\usebox{\mybox}}}

\tikzset{
    triple/.style args={[#1] in [#2] in [#3]}{
        draw,#1,preaction={draw,#2,preaction={draw,#3}}
    }
}   
\tikzset{
  doublebraid/.style={
    triple={[line width=\thickness,white] in [line width=3*\thickness,black] in
    [line width=7*\thickness,white,text=black]}
  }
}

\tikzset{
  functor/.style args={#1}{
    rounded corners,fill=#1,fill opacity=.5,draw opacity=.5
  }
}
\tikzset{
  multimorphism/.style = {regular polygon, regular polygon sides=3,draw,shape border rotate=-90}
}

\tikzset{
  triangle/.style = {draw, regular polygon, regular polygon sides=3, yscale=0.7},
  rotated/.style = {shape border rotate=180},
  dashedstate/.style = {triangle, rotated, dashed},
}
% \def\multimorphism[#1,#2][#3,#4]#5#6#7
%   \path (#1,#2)
%     node [triangle,draw,shape border rotate=-90,inner sep=0pt,label=178:$\vdots$] (a) {#5} (#1,#4)
%     node [triangle,draw,shape border rotate=-90,inner sep=0pt,label=178:$\vdots$] (b) {#6} (#3,($1/2*(#4 + #5)$))
%     node [triangle,draw,shape border rotate=-90,label=135:$a_1$,label=230:$a_n$,label=178:$\vdots$] (#7) {$\theta$} (9,2)
%     node [triangle,draw,shape border rotate=-90,inner sep=-25pt,label=178:$\vdots$] (d) {$\qquad\theta\circ(\theta_1,\ldots,\theta_n)$};
% }

\def\multimorphism[#1,#2][#3][#4][#5][#6] {
    \node[multimorphism] (#3) at (#1,#2) {$#3$};
    \draw (#3.north west) to ++(-.3,0) node[left] (#3x1) {#4};
    \draw (#3.south west) to ++(-.3,0) node[left] (#3xn) {#5};
    \draw (#3.east) to ++(.3,0) node[right] (#3y) {#6};
    \node[anchor=center,left of=#3,yshift=1mm,xshift=2mm] {$\vdots$};
  }

\def\twomultimorphism[#1,#2][#3][#4][#5][#6] {
    \node[multimorphism] (#3) at (#1,#2) {$#3$};
    \draw (#3.north west) to ++(-.3,0) node[left] (#3x1) {#4};
    \draw (#3.south west) to ++(-.3,0) node[left] (#3xn) {#5};
    \draw (#3.east) to ++(.3,0) node[right] (#3y) {#6};
  }

\def\narrowtwomultimorphism[#1,#2][#3][#4][#5][#6] {
    \node[multimorphism] (#3) at (#1,#2) {$#3$};
    \draw ([yshift=-2mm]#3.north west) to ++(-.3,0) node[left] (#3x1) {#4};
    \draw ([yshift=2mm]#3.south west) to ++(-.3,0) node[left] (#3xn) {#5};
    \draw (#3.east) to ++(.3,0) node[right] (#3y) {#6};
  }
\def\wire[#1][#2][#3] {
  \draw ($(#1.east)+(.3,0)$) to[out=0,in=180] ($(#2.#3 west)+(-.3,0)$);
}

\def\comb[#1,#2][#3,#4][#5][#6] {
  \setlength\minimummorphismwidth{6mm}
  \node[morphism] (xi1) at (#1,#2) {#5};
  \node[morphism] (xi2) at (#3,#4) {#6};
  \setlength\minimummorphismwidth{0mm}
  \draw ([xshift=2.5pt]xi1.north east) to ([xshift=2.5pt]xi2.south east);
  \node[dashedmorphism] (f) at ($(xi1.north west)!.5!(xi2.south west)$)
  {\phantom{$f$}};
  \draw (xi1.north west) -- (f.south);
  \draw (xi2.south west) -- (f.north);
  \draw (xi1.south) to ++(0,-.5);
  \draw (xi2.north) to ++(0,.5);
}

\newenvironment{dedication}
  {%\clearpage           % we want a new page          %% I commented this
   \thispagestyle{empty}% no header and footer
   \vspace*{\stretch{1}}% some space at the top
   \itshape             % the text is in italics
   \raggedleft          % flush to the right margin
  }
  {\par % end the paragraph
   \vspace{\stretch{3}} % space at bottom is three times that at the top
   \clearpage           % finish off the page
  }


% metadata
\title{Categories for Cryptographic Composability}
\author{Riley Shahar \\ Advised by Angélica Osorno and Adam Groce}
\date{}

% \keywords{TODO}

% \begin{abstract}
%   TODO
% \end{abstract}

\begin{document}
\maketitle

\begin{frame}{Outline}
	\pause
	\begin{itemize}
		\item Cryptographic composability \pause
		\item Why categories? \pause
		\item Towards a categorical theory of cryptography \pause
		\item Open problems
	\end{itemize}
\end{frame}

\begin{frame}{Composition}
	\only<1>{\[\begin{pic}
				\node[morphism,font=\tiny] (f) at (0,0) {$f$};
				\draw (f.south) to ++(0,-.5)  node[right] {$A$};
				\draw (f.north) to ++(0,.5)  node[right] {$B$};
			\end{pic}\]}

	\only<2>{\[\begin{pic}
				\node[morphism,font=\tiny] (f) at (-.5,0) {$f$};
				\draw (f.south) to ++(0,-.5)  node[right] {$A$};
				\draw (f.north) to ++(0,.5)  node[right] {$B$};

				\node[morphism,font=\tiny] (g) at (.5,0) {$g$};
				\draw (g.south) to ++(0,-.5)  node[right] {$B$};
				\draw (g.north) to ++(0,.5)  node[right] {$C$};
			\end{pic}\]}

	\only<3>{\[
			\begin{pic}
				\node[morphism] (f) {$g \circ f$};
				\draw (f.south) to ++(0,-.65) node[right] {$A$};
				\draw (f.north) to ++(0,.65) node[right] {$C$};
			\end{pic}
			\ =\
			\begin{pic}
				\node[morphism] (g) at (0,.75) {$g\vphantom{f}$};
				\node[morphism] (f) at (0,0) {$f$};
				\draw (f.south) to ++(0,-.3) node[right] {$A$};
				\draw (g.south) to  (f.north);
				\draw (g.north) to ++(0,.3) node[right] {$C$};
			\end{pic}
		\]

		\begin{center}
			Sequential (Vertical) Composition
		\end{center}
	}

	\only<4>{\[
			\begin{pic}
				\node[morphism] (f) {$f \otimes g$};
				\draw (f.south) to ++(0,-.65) node[right] {$A \otimes C$};
				\draw (f.north) to ++(0,.65) node[right] {$B \otimes D$};
			\end{pic}
			=\quad
			\begin{pic}
				\node[morphism] (f) at (-.4,0) {$f$};
				\node[morphism] (g) at (.4,0) {$g\vphantom{f}$};
				\draw (f.south) to ++(0,-.65) node[right] {$A$};
				\draw (f.north) to ++(0,.65) node[right] {$B$};
				\draw (g.south) to ++(0,-.65) node[right] {$C$};
				\draw (g.north) to ++(0,.65) node[right] {$D$};
			\end{pic}
		\]

		\begin{center}
			Parallel (Horizontal) Composition
		\end{center}
	}

	\only<5>{\[
			\begin{pic}
				\node[morphism] (f) at (-1,-1) {$f$};
				\draw (f.south) to ++(0,-.65) node[right] {$A$};
				\draw (f.north) to ++(0,1.5) node {};

				\node[morphism] (g) at (0,-1) {$g$};
				\draw (g.south) to ++(0,-.65) node[right] {$B$};
				\draw (g.north) to ++(0,.65) node {};

				\node[morphism] (h) at (1,-1) {$h$};
				\draw (h.south) to ++(0,-.65) node[right] {$C$};
				\draw (h.north) to ++(0,.65) node {};

				\setlength\minimummorphismwidth{13mm}
				\node[morphism] (s) at (0.5, 0) {$s$};
				\draw (s.north) to ++(0,.5) node {};

				\setlength\minimummorphismwidth{17mm}
				\node[morphism] (t) at (-0.25, 1.0) {$t$};
				\draw (t.north) to ++(0,.5) node[right] {$D$};
			\end{pic}
		\]}

	% In the real world, computer systems are huge, complicated combinations of
	% small pieces.\pause

	% In cryptography, we want to prove that such systems are secure.\pause

	% We'd like to avoid re-proving the security of every component every time for
	% every new system.\pause

	% This problem is called \emph{cryptographic composability}.
\end{frame}

\begin{frame}{Cryptography}
	Cryptography is \emph{the mathematics of secure computation.}\pause

	\begin{multicols}{2}
		\[\begin{pic}
				\node[morphism,font=\tiny] (f) at (0,0) {$f$};
				\draw (f.south) to ++(0,-.5)  node[right] {$A$};
				\draw (f.north) to ++(0,.5)  node[right] {$B$};
			\end{pic}\]\pause

		\includegraphics[scale=0.1]{signal.png}
		\includegraphics[scale=0.25]{gpg.jpg}
	\end{multicols}
\end{frame}

\begin{frame}{Cryptographic Composability}
	Say $f$ and $g$ are secure.\pause

	\begin{multicols}{2}
		\begin{center}
			Is $\begin{pic}
					\node[morphism] (g) at (0,.75) {$g\vphantom{f}$};
					\node[morphism] (f) at (0,0) {$f$};
					\draw (f.south) to ++(0,-.3) node[right] {};
					\draw (g.south) to  (f.north);
					\draw (g.north) to ++(0,.3) node[right] {};
				\end{pic}$?
		\end{center}\pause

		\begin{center}
			Is $\begin{pic}
					\node[morphism] (f) at (-.4,0) {$f$};
					\node[morphism] (g) at (.4,0) {$g\vphantom{f}$};
					\draw (f.south) to ++(0,-.65) node[right] {};
					\draw (f.north) to ++(0,.65) node[right] {};
					\draw (g.south) to ++(0,-.65) node[right] {};
					\draw (g.north) to ++(0,.65) node[right] {};
				\end{pic}$?
		\end{center}
	\end{multicols}
\end{frame}

\begin{frame}{Commitment}
	Alice wants to make a binding guess to Bob, but not reveal it yet. \pause
	She wants a \textbf{commitment protocol}. \pause

	We can use this to play rock-paper-scissors: \pause
	\begin{enumerate}
		\item Alice commits $a\in\{R,P,S\}$ \pause
		\item Bob commits $b\in\{R,P,S\}$ \pause
		\item Alice reveals $a$ \pause
		\item Bob reveals $b$
	\end{enumerate}
\end{frame}

% \begin{frame}{Security Definitions}
% 	What does it mean for a commitment protocol to be \emph{secure}? \pause

% 	A good definition should be:\pause
% 	\begin{itemize}
% 		\item \emph{sound:} capture the intuitively important properties; \pause
% 		\item \emph{usable:} relatively straightforward to prove; \pause
% 		\item \emph{composable:} work well with security of other protocols.
% 	\end{itemize}
% \end{frame}

\begin{frame}{Game-based security}
	\pause List the properties we want our protocol to have.\pause

	Secure commitment should be\pause{} \emph{hiding and binding}.\pause

	\begin{block}{The Hiding Game}
		\begin{enumerate}
			\item Bob picks $m_0,m_1$.\pause
			\item Alice commits $m_b$ at random.\pause
			\item Bob guesses $b'$.
		\end{enumerate}
	\end{block}
\end{frame}

\begin{frame}{Game-based composability}
	\begin{columns}
		\column{0.5\textwidth}
		\begin{block}{The Hiding Game}
			\begin{enumerate}
				\item Bob picks two messages
				\item Alice commits one
				\item Bob guesses which one
			\end{enumerate}
		\end{block}
		\column{0.5\textwidth}
		\begin{block}{Rock-Paper-Scissors}
			\begin{enumerate}
				\item Alice commits $a\in\{R,P,S\}$
				\item Bob commits $b\in\{R,P,S\}$
				\item They both reveal
			\end{enumerate}
		\end{block}
	\end{columns}
	\pause\hrulefill

	Many natural commitment protocols suffer from \emph{malleability.}\pause

	It's taken \emph{decades} for a missing property to be
	noticed~\cite{asher-qkd,bhlmo-qkd}!
\end{frame}

\begin{frame}{Simulation-based security}
	\pause
	Compare the protocol to an \emph{ideal world} with a trusted third party.\pause

	% \begin{block}{Ideal Commitment}\pause
	% 	\begin{enumerate}
	% 		\item Alice sends $m$ to the trusted party.\pause
	% 		\item When Alice wants to reveal, the trusted party sends $m$ to
	% 		      Bob.\pause
	% 	\end{enumerate}
	% \end{block}

	Two probability distributions are \emph{computationally indistinguishable} if
	poly-time algorithms cannot tell between them with nontrivial probability.
	\pause

	A protocol is \emph{secure} if it is computationally indistinguishable
	from the ideal world.\pause
\end{frame}

\begin{frame}[fragile]{Simulation-based composability}
	\pause

	\textbf{Theorem}~\cite{micali-rogaway-1992}.
	\emph{Simulation-secure protocols compose securely in sequences of
		polynomial length.}\pause

	\begin{tikzpicture}
	\end{tikzpicture}

	However,~\cite{goldreich-krawczyk-1996} gave a protocol for zero-knowledge
	proof that's simulation secure, but doesn't compose in parallel.
\end{frame}

\begin{frame}[fragile]{Universal Composability}
	\pause

	\textbf{Theorem}~\cite{canetti-2000a}.
	\emph{UC-secure protocols compose securely in parallel sequences of polynomial
		length or width.}\pause

	UC was revised in
	~\cite{canetti-2000a,canetti-2000b,canetti-2001,canetti-2005a,canetti-2005b,canetti-2013a,canetti-2013b,canetti-2018,canetti-2020}.\pause

	The proofs\dots\begin{itemize}
		\item are wildly dependent on small technical details;\pause
		\item leave artifacts in the protocol;\pause
		\item are very hard to trust.
	\end{itemize}
\end{frame}

\begin{frame}[fragile]
	\parbox{5cm}{Cryptography is in need of an elegant mathematical theory abstracting
		composability of computational processes\dots}\pause
	\vfill
	\hfill\parbox{5cm}{\dots \textbf{category theory} is an excellent candidate
		for such a theory.}
\end{frame}

\begin{frame}{Categories and Programming Languages}
	\begin{multicols}{2}
		In programming language theory:\pause
		\begin{itemize}
			\item \only<2->{objects are types;}
			      \only<3->{\item morphisms are programs.}
		\end{itemize}

		\[
			\begin{pic}
				\only<3->{\node[morphism,font=\tiny] (f) at (0,0) {$len$};}
				\only<2->{\draw (f.south) to ++(0,-.5)  node[right] {\texttt{str}};
					\draw (f.north) to ++(0,.5)  node[right] {\texttt{nat}};}
			\end{pic}
		\]
	\end{multicols}

	\only<4->{(We've already been doing category theory!)}
\end{frame}

\begin{frame}{Composing Programs}
	\begin{multicols}{2}
		Consider programs \hspace{0.5em} \[
			\begin{pic}
				\node[morphism,font=\tiny] (f) at (0,0) {$f$};
				\draw (f.south) to ++(0,-.5)  node[right] {$A$};
				\draw (f.north) to ++(0,.5)  node[right] {$B$};
			\end{pic} \text{ and }\hspace{0.5em}
			\begin{pic}
				\node[morphism,font=\tiny] (g) at (0,0) {$g$};
				\draw (g.south) to ++(0,-.5)  node[right] {$B$};
				\draw (g.north) to ++(0,.5)  node[right] {$C$};
			\end{pic}.\]\pause

		We can \emph{always} make a program\[
			\begin{pic}
				\node[morphism] (g) at (0,.75) {$g\vphantom{f}$};
				\node[morphism] (f) at (0,0) {$f$};
				\draw (f.south) to ++(0,-.3) node[right] {$A$};
				\draw (g.south) to  (f.north);
				\draw (g.north) to ++(0,.3) node[right] {$C$};
			\end{pic}.\]
	\end{multicols}

	Composition ``works'' in PL theory.
\end{frame}

\begin{frame}{Categorical Cryptography: The Idea}
	Make a relation $\approx$ between morphisms, roughly like
	indistinguishability.\pause

	\[ \begin{pic}[yscale=-1]
			\node[dot,fill=white] (d) {};
			\draw (d) to +(0,0.5) node [left] {$A$};
			\draw (d) to[out=0,in=90] +(0.5,-0.5) node[right] {$B$};
			\draw (d) to[out=180,in=90] +(-0.5,-0.5) node[left] {$E$};
		\end{pic}
		\otimes
		\begin{pic}
			\node[dot,fill=white] (d) {};
			\draw (d) to +(0,-0.5) node[state,scale=0.75] {\normalsize$\$$};
			\draw (d) to[out=180,in=-90] +(-0.5,0.5) node[left] {$A$};
			\draw (d) to[out=0,in=-90] +(0.5,0.5) node[right] {$B$};
		\end{pic}
		\xrightarrow{\text{One-Time Pad}}
		\begin{pic} \draw (0,0) node[left] {$A$} to (0,1) node[right] {$B$}; \end{pic}
		\otimes
		\hspace{0.5em}
		\begin{pic} \draw (0,0) node[state,scale=0.75] {\normalsize$\$$} to (0,0.5) node[right] {$E$}; \end{pic}
	\]
\end{frame}

\begin{frame}[fragile]{Categorical Composable Cryptography}
	As far as I know, there is only one published paper using category theory for
	composability in cryptography~\cite{broadbent-karvonen-2022}.\pause

	% \only<2-2>{
	\begin{equation*}
		\begin{tikzpicture}[scale=0.5, baseline=(c)]]
			\path coordinate[dot, fill=white] (c)
			+(-1,1) coordinate[label=above:$E$] (tl)
			+(1,1) coordinate[label=above:$B$] (tr)
			+(0,-1) coordinate[label=below:$A$] (bot)
			(tl) ++(-0.5,0) coordinate (tll)
			(bot) ++(1.5,0) coordinate (br);
			\draw (bot) -- (c) to[out=0, in=-90] (tr);
			\draw (c.center) to[out=180, in=-90] (tl);
		\end{tikzpicture}
		\otimes
		\begin{tikzpicture}[scale=0.5, baseline=(c)]]
			\path coordinate[dot,fill=white] (c)
			+(-1,1) coordinate[label=above:$A$] (tl)
			+(1,1) coordinate[label=above:$B$] (tr)
			+(0,-1) coordinate[label=below:$\$$] (bot)
			(tl) ++(-0.5,0) coordinate (tll)
			(bot) ++(1.5,0) coordinate (br);
			\draw (bot) -- (c) to[out=0, in=-90] (tr);
			\node[state=-90, anchor=north] at (bot) {};
			\draw (c.center) to[out=180, in=-90] (tl);
		\end{tikzpicture}
		\xrightarrow{\text{One-Time Pad}}
		\begin{tikzpicture}[scale=0.5, baseline=($(bot)!0.5!(top)$)]]
			\path
			+(0,-1) coordinate[label=below:$A$] (bot)
			+(0,1) coordinate[label=above:$B$] (top);
			\draw (bot) -- (top);
			% \draw (c.center) to[out=180, in=-90] (tl);
		\end{tikzpicture}
		\otimes
		\begin{tikzpicture}[scale=0.5,baseline=($(bot)!0.5!(top)$)]
			\path
			+(0,-0.5) coordinate[label=below:$\$$] (bot)
			+(0,0.5) coordinate[label=above:$E$] (top);
			\draw (bot) -- (top);
			\node[state=-90, anchor=north] at (bot) {};
		\end{tikzpicture}
	\end{equation*}
	% }

	\pause
	The moral:\begin{itemize}\pause
		\item objects are resources, like channels or keys;\pause
		\item morphisms are ``protocols with holes'';\pause
		\item composition ``plugs in the holes''.\pause
	\end{itemize}

	This is the standard paradigm in programming language theory, transported to
	cryptographic language.
\end{frame}

\begin{frame}{Attack Models}
	In CCC, adversaries are constrained by \emph{attack models}.\pause

	An attack model $\cA$ assigns, to each morphism $f$, a collection of morphisms
	$\cA(f)$ satisfying some axioms. If the adversary is ``supposed'' to do $f$,
	then they can instead do anything in $\cA(f)$.\pause

	\textbf{Open Question 1:} \emph{Can the axioms be formulated as functoriality plus
		some conditions? If $\cA$ is a functor, what should its codomain be?}\pause

	\textbf{Open Question 2:} \emph{How broad is the definition of an attack
		model? Does it capture enough of modern cryptography?}
\end{frame}

\begin{frame}{Computational Indistinguishability}
	Composition should only work polynomially many times.\pause{} In fact,
	computational indistinguishability is not even an equivalence relation.\pause

	B\&K work around this by artificially limiting the universe size.\pause

	\textbf{Open Question 3:} \emph{Is there a natural model of computational
		indistinguishability in nice symmetric monoidal categories?}\pause

	B\&K propose attaching an extended metric to the category, interpreted as
	the ``computational distance'' between two morphisms.\pause

	We need some way to deal with asymptotic behavior. Our current idea is to
	value the metric in $\RR^{\NN}$.
\end{frame}

\begin{frame}{Enrichment}
	Many of the constructions of B\&K enhance the category with extra
	structure.\pause{}

	\textbf{Open Question 4:} \emph{Do their composition theorems extend to the context of
		enriched category theory?\pause{} If so, can the extra structure they need be framed
		as some kind of enrichment?}
\end{frame}

\begin{frame}{Conclusion}
	CCC is very important progress, but there's a lot to be done.\pause{} CCC
	largely staples cryptographic notions on top of a category, rather than working
	with the machinery. We'd like to try to improve on this.\pause

	Once we have a more natural model, there will be lots of cool applications:\pause
	\begin{itemize}
		\item Can we incorporate categorical notions from game theory, programming
		      languages, etc. into cryptography?\pause
		\item What does the presence of various categorical structure ((co)limits,
		      monads, etc.) say cryptographically?
	\end{itemize}
\end{frame}
\end{document}
